% fig07-bitmask.tex -- Figure 2.2: bitmask <-> subset for n = 3.
% One row per mask: the integer, its bits, the subset it encodes.
\begin{tikzpicture}[font=\small, every node/.style={inner sep=0pt}]
\node[draw, text width=0.56\linewidth, inner sep=8pt, font=\footnotesize,
      align=flush left] (tab) at (0,0) {%
  \begin{tabular}{cccl}
  \toprule
  \textbf{mask} & \textbf{bits} $b_2 b_1 b_0$ & \textbf{subset} & \textbf{loop order}\\
  \midrule
  0 & 000 & $\{\}$        & first\\
  1 & 001 & $\{1\}$       &\\
  2 & 010 & $\{2\}$       &\\
  3 & 011 & $\{1,2\}$     &\\
  4 & 100 & $\{3\}$       &\\
  5 & 101 & $\{1,3\}$     &\\
  6 & 110 & $\{2,3\}$     &\\
  7 & 111 & $\{1,2,3\}$   & last\\
  \bottomrule
  \end{tabular}};
\draw[{stealth}-, thick, accent] ($(tab.north east)+(-2.0cm,-0.15)$)
  -- ++(0.5,0.55) node[right, font=\footnotesize, text=accent,
                      text width=3.4cm, align=flush left]
  {bit $k$ set $\iff$ element $k{+}1$ chosen};
\draw[{stealth}-, thick, accent] ($(tab.east)+(0,-0.55)$)
  -- ++(0.5,-0.3) node[right, font=\footnotesize, text=accent,
                      text width=4.0cm, align=flush left]
  {$2^n$ masks cover every subset exactly once};
\end{tikzpicture}
